% !TEX program = lualatex
% =============================================================
% OnePage LG — Before / After / Result
%
%   ESTE ARQUIVO E SO CONTEUDO. Preencha os campos abaixo e compile.
%   Layout, cores, fontes e coordenadas ficam em onepagelg.sty.
%
%   Os itens repetiveis (\ponto, \destaque, \melhoria, \afirmacao, \efeito,
%   \area, \graficobarra) se empilham sozinhos: acrescente ou remova linhas
%   a vontade, sem mexer em coordenada nenhuma.
%
%   Os icones ficam em img/. Trocar o desenho e trocar o nome do arquivo.
%
%   Compilar: lualatex onepage-lg.tex   (2x — a pagina usa remember picture)
% =============================================================
\documentclass{article}
\usepackage{onepagelg}
\begin{document}

% ======================= CABECALHO E RODAPE =======================
\avisotopo{LGE Internal Use Only}
\titulo{1. AI Deep Dive}
\subtitulo{Framework to automate business processes through AI tools and promote LGEAZ as AI Driven Corporation.}
\colunas{Before}{After}{Result}
\rodape{| \ C\,O\,N\,F\,I\,D\,E\,N\,T\,I\,A\,L\ |}{0 / 7}

% ======================= COLUNA 1 — BEFORE =======================
\tituloantes{Background / Issue}
\textoantes{High dependence on repetitive manual tasks, consuming many hours and
  reducing productivity and operational efficiency across all departments.}

% \ponto[altura do icone]{icone}{texto}
\ponto[1.26]{ico-repetitivo.png}{75\% of daily and weekly activities, consuming many hours in each department.}
\ponto[1.14]{ico-sobrecarga.png}{SCM department overload.}
\ponto[0.95]{ico-planilhas.png}{Data mostly base on spreadsheets}

\alerta{ico-seguranca.png}{Information security risks}

% quadro de baixo: titulo, o fluxo "de ... para ..." e os numeros de impacto
\tituloquadro{Current situation}
\fluxo{40K hours saved}{15\% of total digitization}
\destaque{ico-pessoas.png}{138 Overworked Employees}
\destaque{ico-oportunidades.png}{95\% of opportunities not covered by current technologies}

% ======================= COLUNA 2 — AFTER =======================
\titulodepois{Expected Improvement}
\textodepois{With the implementation of technologies capable of overcoming the
  challenge of digital transformation, interconnected processes managed by
  AI-based automated solutions are expected.}

\melhoria{40\% increase in the level of automation.}
\melhoria{Monitor impacts of each solution in real time.}

\abas{AI DEEP DIVE}{AI CORE}

% \area{angulo}{rotulo} — o rotulo se posiciona sozinho pelo angulo
\area{90}{SCM}
\area{45}{Material}
\area{0}{Production}
\area{-45}{Planning}
\area{270}{Procurement}
\area{225}{QA}
\area{180}{RH}
\area{135}{Tax}

\afirmacao{Project team capable of developing intelligent tools and applied research.}
\afirmacao{Developers directly connected to business areas.}
\afirmacao{Able to understand key value process and generate solutions.}

\ilustracoes{robo-analista.png}{troca.png}{lg-ia.png}

% ======================= COLUNA 3 — RESULT =======================
\tituloresultado{Improvement Effect}

\efeito{Empowering digital transforming}
\efeito{Promote governance and performance measurement}
\efeito{Optimize process flows across multiple areas}

% \graficobarra{categoria}{rotulo do valor}{altura em cm}
\graficoeixo{FTE}
\graficobarra{Before}{40K}{1.04}
\graficobarra{After}{150K}{2.42}
\graficonota{Multiple areas}

\impacto{ico-areas.png}{Department Impacted}{Multiple Areas}
\resultado{FTE (Expected)}{150K Hours}

% =============================================================
\onepage
\end{document}
